% =====================================================================
% Conference paper in ifacconf format, targeting ADCHEM 2027 (Hong Kong)
% or similar IFAC venue. 6-page target.
%
% Front matter adapted from elsarticle version; body text unchanged.
%
% Numbers policy: all from verified follow-up runs. Entries marked
% TODO-REGEN need regeneration before submission; TODO mark editorial
% placeholders.
% =====================================================================
\documentclass[twocolumn]{ifacconf}

\usepackage{amsmath,amssymb,amsthm}
\usepackage{graphicx}
\usepackage{booktabs}

\newtheorem{proposition}{Proposition}
\newtheorem{corollary}{Corollary}
\newtheorem{lemma}{Lemma}
\newtheorem{remark}{Remark}

\newcommand{\Nbar}{\bar{N}}
\newcommand{\Neps}{N_{\varepsilon}}

\title{Limit Cycles and Regularization in Step-Response Phase-Portrait PID Tuning}

\author{Author One$^{a,}$\thanks{Corresponding author: TODO-EMAIL}, 
        Author Two$^b$}

\address{
$^a$TODO Affiliation One, City, Country. \\
$^b$TODO Affiliation Two, City, Country.
}

\begin{document}

\maketitle

\begin{abstract}
In a companion paper we introduced SPIN, a model-free PID tuning method
that inspects three phase portraits of the closed-loop step response,
counts a turn index in each, and adjusts the gains by a fixed table of
four moves. The method works, but the companion paper does not say what
the iteration converges \emph{to}. This paper answers that question.
First, we show that in log-gain coordinates the four moves admit a unique
vanishing convex combination, so the only possible resting target is the
\emph{triple point}: the gain vector at which all three turn indices sit
exactly on their thresholds. When the triple point falls outside the gain
box, the target degenerates to a complementarity point, which explains a
boundary run reported in the companion paper. Second, we show that the
published turn index is a discontinuous function of the gains, so the
triple-point equation generally has no solution and diminishing-step runs
freeze. Third, we repair this with an $\varepsilon$-regularized winding
integral, after which diminishing-step runs converge to the triple point.
Finally, we quantify the limit cycles of constant-step runs and compare
step-size schedules; a simple reversal-triggered halving schedule reaches
$10^{-4}$ accuracy in under a hundred iterations.
\end{abstract}

\keywords{PID control, model-free tuning, winding number, switched systems,
limit cycles}

% =====================================================================
\section{Introduction}
\label{sec:intro}

Most PID tuning methods fit a model first and compute gains second
\cite{ZN1942,AstromHagglund2004}. In a companion paper \cite{Paper1} we took
a different route: tune directly from what the step response \emph{looks like}.
The method, SPIN (Step-response Phase-portrait INspection), draws three phase
portraits from one closed-loop step response, counts how many times each curve
turns around the origin --- the \emph{turn index} $N_k$ of band $k$ --- and
compares each count to a fixed threshold. A table of four moves, the
\emph{triangular rule}, then scales the gains up or down. No plant model is
identified at any point.

The companion paper demonstrates that the rule detunes aggressive gains,
recovers from failed loops, and lands in a feasible region on a battery of
standard plants. What it does not do is characterize the iteration's
asymptotic behavior. Three questions were left open. What point, if any, does
the rule aim at? Why do runs with a shrinking step size freeze at points that
depend on the schedule rather than on the plant? And which step-size schedules
actually converge?

This paper answers all three. The answers turn out to be connected. The rule
aims at a single well-defined target, the \emph{triple point}, at which all
three turn indices sit exactly on their thresholds (Section~\ref{sec:triple}).
Diminishing-step runs freeze because the published turn index is discontinuous
in the gains, so the triple-point equation generally has no solution
(Section~\ref{sec:obstruction}). Replacing the hard count with an
$\varepsilon$-regularized winding integral removes the discontinuity, and with
it the freezing (Section~\ref{sec:regularized}). With a continuous index in
hand, the limit cycles of constant-step runs and the convergence of
diminishing-step schedules can be measured cleanly (Section~\ref{sec:schedules}).
Section~\ref{sec:structure} examines the local structure at the target and
identifies the route a convergence proof must take.

It is worth stating what kind of object we are analyzing. Iterative model-free
tuning is not new: iterative feedback tuning estimates the gradient of a
performance criterion from closed-loop experiments and converges, with an
appropriately shrinking step size, to a local minimum of that criterion
\cite{Hjalmarsson1998}; extremum seeking does the same with sinusoidal probing
\cite{Killingsworth2006}. The triangular rule is different in kind. It is a
\emph{sign-based} update: each iteration compares three integers (later: three
real numbers) to three thresholds and picks one of four fixed moves. There is
no cost function and no gradient, so the limit object is not a stationary point
of anything. It is a point where a switching law balances --- and the right
mathematical frame is that of switched systems and their averaged dynamics, not
descent.

% =====================================================================
\section{The Rule as a Switched System}
\label{sec:rule}

We recall only what is needed; the companion paper \cite{Paper1} defines the
portraits, the index, and the rule in full. One closed-loop step experiment
yields three curves (band $0$, $1$, $2$), and for each curve a turn index
$N_k \ge 0$ measuring how far the curve winds around the origin of its
portrait. Each band has a fixed threshold:
\begin{equation}
\Nbar = (\Nbar_0, \Nbar_1, \Nbar_2) = (0.5,\; 0.75,\; 1.0).
\label{eq:thresholds}
\end{equation}

The triangular rule reads the three comparisons $N_k \gtrless \Nbar_k$ and
applies one of four moves to the gain vector $K = (K_p, K_i, K_d)$: if all
three bands are below threshold, all gains are scaled up; otherwise the lowest
offending band is cut by scaling gains down in a band-specific pattern. Every
move multiplies each gain by $(1+\beta)^{\pm 1}$ or leaves it alone, where
$\beta > 0$ is the step size. Gains are clipped to a box, here $[0.01, 10]$.

In log coordinates $\theta = (\log K_p, \log K_i, \log K_d)$, the four moves
become translations by $h = \log(1+\beta)$:
\begin{equation}
\begin{aligned}
e   &= (+1, +1, +1) &\text{(expand),}\\
d_0 &= (-1, \phantom{+}0, \phantom{+}0) &\text{(cut band 0),}\\
d_1 &= (+1, -1, \phantom{+}0) &\text{(cut band 1),}\\
d_2 &= (+1, +1, -1) &\text{(cut band 2),}
\end{aligned}
\label{eq:moves}
\end{equation}
each scaled by $h$. The iteration is a switched dynamical system: a partition
of gain space into four regions (determined by the signs of $N_k(\theta) -
\Nbar_k$), with a constant velocity assigned to each region. Two properties
drive everything that follows. First, velocities are constant within each
region, so convergence requires the state to reach where regions meet and moves
balance. Second, switching surfaces are the level sets $N_k(\theta) = \Nbar_k$,
so their regularity --- or lack of it --- decides whether such a point exists.

% =====================================================================
\section{The Target: a Triple Point}
\label{sec:triple}

Suppose a trajectory stays bounded and does not drift. Over a long window, each
move $v \in \{e, d_0, d_1, d_2\}$ is applied some fraction $w_v \ge 0$ of the
time, with $\sum_v w_v = 1$. Non-drifting requires the average displacement
$h \sum_v w_v v$ to vanish.

\begin{proposition}[Unique Balance]
\label{prop:balance}
The only convex combination of the moves \eqref{eq:moves} that vanishes is
\begin{equation}
\tfrac{1}{8}\, e + \tfrac{1}{2}\, d_0 + \tfrac{1}{4}\, d_1
+ \tfrac{1}{8}\, d_2 = 0.
\label{eq:weights}
\end{equation}
\end{proposition}

All four weights are strictly positive. This has a sharp consequence: a bounded,
non-drifting trajectory must keep visiting all four regions, possible only
where all three switching surfaces meet.

\begin{corollary}[Triple Point]
\label{cor:triple}
The only interior resting target is a point $\theta^*$ with
\begin{equation}
N_k(\theta^*) = \Nbar_k, \qquad k = 0, 1, 2,
\label{eq:triplepoint}
\end{equation}
approached through a period-8 orbit with move frequencies
$(\tfrac18, \tfrac12, \tfrac14, \tfrac18)$.
\end{corollary}

This is testable. On plants P1, P3, P4 of the companion battery, long
constant-step runs settle into a periodic orbit, and the measured duty cycle
matches $(\tfrac18, \tfrac12, \tfrac14, \tfrac18)$ to three decimal places, at
both $\beta = 0.10$ and $\beta = 0.03$. The terminal orbit is a period-8 cycle
in which the moves appear in a fixed ``ruler'' pattern.

\begin{table}[t]
\small
\centering
\caption{Measured duty cycle vs.\ prediction.
% TODO-REGEN: values from code/asymptotics/duty_cycle.py; matched
% to three decimals on P1, P3, P4 at both step sizes.
}
\label{tab:duty}
\begin{tabular}{lcccc}
\toprule
 & $w_e$ & $w_{d_0}$ & $w_{d_1}$ & $w_{d_2}$ \\
\midrule
Predicted & $0.125$ & $0.500$ & $0.250$ & $0.125$ \\
P1, P3, P4 & \multicolumn{4}{c}{TODO-REGEN} \\
\bottomrule
\end{tabular}
\end{table}

\paragraph{Boundary Case} If $\theta^*$ falls outside the gain box, condition
\eqref{eq:triplepoint} cannot be met. The target degenerates: for each band,
either the index sits on its threshold \emph{or} the gain sits on a bound --- a
complementarity point. The companion paper's P2 run ends with $K_d$ on the
$0.01$ floor and indices $N = (0.11,\, 1.04,\, 1.14)$. Two indices sit at
thresholds; the constraint is replaced by an active gain bound. What the
companion paper calls ``the rule removing derivative action on a delay-dominant
plant'' is exactly this complementarity target.

% =====================================================================
\section{The Obstruction: a Discontinuous Count}
\label{sec:obstruction}

Corollary~\ref{cor:triple} assumes equation \eqref{eq:triplepoint} has a
solution. With the published turn index, it generally does not. The index is a
hard count: the portrait curve is followed, winding is accumulated, and
accumulation is truncated when the curve enters a small disc of radius
$\varepsilon$ about the origin for the last time. The truncation is the
problem. The time of \emph{last entry} into the disc is not a continuous
functional of the curve: a small deformation can make the curve graze the disc
boundary and either gain or lose a whole excursion.

On the battery plants, a gain change of $0.05\,\%$ --- far below any step the
rule takes --- moves one turn index by exactly $1.000$. Lengthening the window
does not help: the jump is caused by the truncation rule. The consequence for
the asymptotics is immediate. The switching surfaces $N_k(\theta) = \Nbar_k$
are, in general, not surfaces: the index jumps across them in integer steps, so
the triple-point equation \eqref{eq:triplepoint} often has no solution, and a
run whose step size shrinks to zero must freeze at whatever point remains ---
a point depending on the schedule, not the plant.

For a constant step this discontinuity is harmless: the rule only reads the
\emph{sign} of $N_k - \Nbar_k$, and a jump merely relocates the switching
surface slightly. This is why the companion paper's results hold. The
discontinuity only becomes fatal when one asks the iteration to converge to a
point.

% =====================================================================
\section{The Fix: a Regularized Turn Index}
\label{sec:regularized}

The repair is to stop truncating and start weighting. For a closed curve
$(x(t), y(t))$ not passing through the origin, the winding number is
$\frac{1}{2\pi} \oint \frac{x\,dy - y\,dx}{x^2 + y^2}$. We damp instead of
truncate:
\begin{equation}
\Neps \;=\; \frac{1}{2\pi} \oint
\frac{x\,dy - y\,dx}{\,x^2 + y^2 + \varepsilon^2\,}.
\label{eq:neps}
\end{equation}

Writing $\rho^2 = x^2 + y^2$, the integrand is the ordinary winding integrand
multiplied by $\rho^2 / (\rho^2 + \varepsilon^2)$: turning far from the origin
counts fully, turning at radius $\varepsilon$ counts half, and turning deep
inside counts almost nothing. The disc's sharp boundary is gone, and with it
the need to decide when the curve entered the disc ``for the last time''.

\emph{Continuity:} The integrand is smooth everywhere, so $\Neps$ is continuous
in the gains. On the $0.05\,\%$ perturbation, the change is $7 \times 10^{-4}$
where the hard count jumps by $1.000$.

\emph{Window independence:} The tail of a settled response contributes radius
$\rho \to 0$, where the weight vanishes quadratically. The value becomes
independent of the window once the response settles, and the $\delta$-guard is
unnecessary.

\emph{Interchangeability:} Away from switching surfaces, $\Neps$ and the hard
count agree to within a fraction vanishing as $\varepsilon \to 0$, so every
sign comparison the published rule makes is unchanged. All companion-paper
results hold with $\Neps$ in place of the hard count.

With $\Neps$, the switching surfaces are genuine level sets, equation
\eqref{eq:triplepoint} has a solution, and diminishing-step runs stop freezing.
On the three interior-target plants, runs with vanishing steps now converge to
the triple point with final log-gain errors
\begin{equation*}
\text{P1: } 0.010, \quad \text{P3: } 0.017, \quad \text{P4: } 0.041.
\end{equation*}

% =====================================================================
\section{Step-Size Schedules and Limit Cycles}
\label{sec:schedules}

All runs start far from the target --- gains detuned by $16\times$ --- and use
the regularized index. We compare constant step against diminishing schedules
and an adaptive rule. Table~\ref{tab:schedules} summarizes the outcome; errors
are final log-gain distances to $\theta^*$.

\begin{table}[t]
\small
\centering
\caption{Final log-gain error for different step-size schedules.
% TODO-REGEN: from code/asymptotics/schedule_comparison.py; rows may
% be per-plant or averaged, confirm which and add to reproduce script.
}
\label{tab:schedules}
\begin{tabular}{lll}
\toprule
Schedule & Final error & Behaviour \\
\midrule
$\beta$ const.           & $0.074$ & limit cycle \\
$\beta_n \propto n^{-1/2}$ & $0.010$ & converges \\
$\beta_n \propto n^{-3/4}$ & $0.001$ & converges \\
$\beta_n \propto n^{-1}$   & $0.197$ & freezes \\
reversal-halving         & $10^{-4}$ & $36$--$98$ iters \\
\bottomrule
\end{tabular}
\end{table}

\paragraph{Constant Step} A constant $\beta$ cannot converge: near $\theta^*$ the
state orbits with radius scaling $h = \log(1+\beta)$. The error $0.074$ is this
orbit radius. For plant use, gains within a few percent are indistinguishable.

\paragraph{Diminishing Steps} Schedules $n^{-1/2}$ and $n^{-3/4}$ converge; the
classical $n^{-1}$ freezes short at error $0.197$ for a travel-budget reason:

\begin{lemma}[Logarithmic Travel]
\label{lem:travel}
Under $\beta_n = c/n$, total distance in $n$ steps is $\sum_{m \le n}
\log(1 + c/m) = O(\log n)$. A far start at distance $D$ cannot be reached
before $n \approx e^{D/c}$ iterations.
\end{lemma}

With a far start ($D \approx \log 16$ per gain), the required $n$ is
astronomically large. This mirrors the summability boundary in stochastic
approximation, met here from the travel side rather than noise.

\paragraph{Reversal-Triggered Halving} Hold $\beta$ constant until every band
has reversed its move direction at least once, then halve $\beta$ and repeat.
Reversals indicate the trajectory has begun to orbit rather than travel, so the
step is no longer limiting and can be cut. This schedule reaches error $10^{-4}$
in $36$--$98$ iterations. The mechanism is a global cousin of the Rprop
update with reversal-triggered step adaptation \cite{Riedmiller1993}, arising
naturally here because the triangular rule is sign-based.

% =====================================================================
\section{Local Structure and Proof Route}
\label{sec:structure}

Central-difference Jacobians $\partial N_j / \partial \theta_i$ at $\theta^*$ are
nonsingular on all interior-target plants (determinants $1.2$, $5.9$, $4.3$),
so $\theta^*$ is locally unique.
% TODO-REGEN: add units/normalization to reproduce script.

The Jacobians are decidedly \emph{not} diagonal or diagonally dominant. On P1,
proportional gain moves band-2 index harder than derivative gain does; on P4,
the own-gain entry of band 1 is negative. The intuitive picture --- each move
mainly fixes its own band --- is false at the target: per-move descent fails on
two of three plants. Yet all runs converge.

The resolution is that convergence is a property of the \emph{averaged}
dynamics, not individual moves. Near $\theta^*$ the iteration chatters between
regions, and its mean motion follows the convex combination of region
velocities weighted by local residence times --- the Filippov regularization
\cite{Filippov1988}. The duty-cycle measurements of Section~\ref{sec:triple}
are direct evidence: observed residence times are exactly the balance weights
\eqref{eq:weights}. A convergence proof should establish stability of $\theta^*$
for the Filippov dynamics and conclude for the discrete iteration by standard
averaging. Per-band acceleration (independent step sizes or momentum) assumes
diagonal structure that is absent and should be expected to misbehave; momentum
on the full move vector is the variant worth trying.

% =====================================================================
\section{Conclusions}

The triangular rule is a switched system whose only interior target is the
triple point: all three indices exactly on threshold, approached through a
period-8 orbit with fixed move frequencies. When the triple point leaves the
gain box, the target is a complementarity point (companion paper's P2 run). The
hard-truncated count is discontinuous and cannot support point convergence; the
$\varepsilon$-regularized winding integral restores continuity at no cost. With
it, diminishing-step runs converge, $n^{-1}$ fails for a travel-budget reason,
and reversal-triggered halving reaches $10^{-4}$ in under a hundred iterations.
The convergence proof itself runs through averaged Filippov dynamics and is left
open, together with the momentum variant it suggests.

% =====================================================================
\begin{thebibliography}{10}

\bibitem{Paper1}
Authors, Model-free PID tuning by step-response inspection, submitted to
J. Process Control, 2026.

\bibitem{ZN1942}
J.~G. Ziegler, N.~B. Nichols, Optimum settings for automatic controllers,
Trans. ASME 64 (1942) 759--768.

\bibitem{AstromHagglund2004}
K.~J. {\AA}str{\"o}m, T. H{\"a}gglund, Revisiting the Ziegler--Nichols
step response method for PID control, J. Process Control 14 (2004) 635--650.

\bibitem{Hjalmarsson1998}
H. Hjalmarsson, M. Gevers, S. Gunnarsson, O. Lequin, Iterative feedback
tuning: theory and applications, IEEE Control Syst. Mag. 18 (4) (1998)
26--41.

\bibitem{Killingsworth2006}
N.~J. Killingsworth, M. Krsti{\'c}, PID tuning using extremum seeking,
IEEE Control Syst. Mag. 26 (1) (2006) 70--79.

\bibitem{Riedmiller1993}
M. Riedmiller, H. Braun, A direct adaptive method for faster backpropagation
learning: the RPROP algorithm, in: Proc. IEEE Int. Conf. Neural Networks,
1993, pp. 586--591.

\bibitem{Filippov1988}
A.~F. Filippov, Differential Equations with Discontinuous Righthand Sides,
Kluwer, 1988.

\end{thebibliography}

\end{document}
